% ECONOMICS_PROBLEM: {"id": "D44-196", "title": "Computing median prices for XOS valuations", "jel_code": "D44", "source_id": "OP-0176"}
% !TeX program = xelatex
\documentclass[11pt,a4paper]{article}
\usepackage[margin=23mm]{geometry}
\usepackage{fontspec}
\setmainfont{Latin Modern Roman}
\usepackage{amsmath,amssymb,mathtools}
\usepackage{xcolor,enumitem,booktabs,longtable,array,xurl,hyperref}
\definecolor{muted}{HTML}{53636C}
\hypersetup{colorlinks=true,urlcolor=blue,linkcolor=blue}
\setlength{\parindent}{0pt}
\setlength{\parskip}{5pt}
\newcommand{\R}{\mathbb R}
\newcommand{\E}{\mathbb E}
\newcommand{\N}{\mathbb N}
\newcommand{\ind}{\mathbf 1}
\newcommand{\Var}{\operatorname{Var}}
\newcommand{\Cov}{\operatorname{Cov}}
\newcommand{\supp}{\operatorname{supp}}
\DeclareMathOperator*{\argmin}{arg\,min}
\DeclareMathOperator*{\argmax}{arg\,max}
\newcommand{\entrymeta}[3]{{\sffamily\small\color{muted}\textbf{\texttt{#1}}\quad|\quad #2\par #3\par}}
\newcommand{\Problemref}[1]{\texttt{#1}}
\newcommand{\JELref}[2]{\texttt{#2} (JEL \texttt{#1})}
\begin{document}
\section*{Computing median prices for XOS valuations}
\textbf{Problem ID:} \texttt{D44-196}\quad \textbf{JEL:} \texttt{D44}\par
% BEGIN ECONOMICS STATEMENT
\label{problem:OP-0176}
% GAME_THEORY_PROBLEM: {"local_id": "GT-046", "original_number": 46, "kind": "P", "entry_kind": "statement_only", "published": false, "review_required": true, "category": "game_theory"}
\label{game:GT-046}
{\sffamily\small\color{muted}
\textbf{GT-046} | Online computational mechanisms | \textbf{P}: source-explicit computational question.
\par}
\subsection*{Definitions and mathematical statement}
Let $D=\bigotimes_{i=1}^nD_i$ be a generic distribution over XOS valuations on $M=[m]$. XOS means $v_i(S)=\max_{\ell}\sum_{g\in S}a_{i\ell g}$ for nonnegative coefficients. Generic means that for each fixed price vector and available set, the utility-maximizing bundle is unique almost surely. Buyers arrive in fixed order and choose demanded subsets of remaining goods. Let $\pi_g(p)$ be the probability that $g$ is sold at prices $p$. A nonnegative price vector $p$ is $\alpha$-median if
\[
 \pi_g(p)\leq1-\alpha,\qquad
 p_g>0\ \Longrightarrow\ \pi_g(p)\geq\alpha
 \quad(g\in M).
\]
The computational question is whether, from independent samples of $D$, one can output such a vector for $\alpha=1/2-\varepsilon$, with probability at least $1-\delta$, in time and sample size polynomial in $n,m,1/\varepsilon,\log(1/\delta)$ and the numerical input lengths. Valuations are accessed through demand and value oracles, and their oracle-call costs are counted; zero-price items need no lower sale-probability bound. The input accuracy parameters satisfy $0<\varepsilon<1/2$ and $0<\delta<1$.
\subsection*{Short English statement}
Extend efficient median-price learning from unit-demand buyers to XOS buyers.
\subsection*{Variants and scope boundaries}
Existence and polynomial sample complexity do not establish polynomial computation. An algorithm may randomize while returning a single price vector. With atoms, tie-breaking must be incorporated using the source's nongeneric extension. A mixture satisfying merely average sale bounds is a different output target.
\subsection*{Source relationship and status}
Section 4.5 explicitly leaves the XOS extension of Theorem 4.9 open. The formal condition here is Definition 4.2, and the unit-demand proof constructs a single vector. The oracle/encoding conventions above make computational efficiency meaningful; they do not imply that demand-oracle access is equivalent to a full explicit valuation table.
\subsection*{References}
\begingroup\footnotesize\raggedright
\textbf{[S1]} Paul D\"utting, Thomas Kesselheim, Brendan Lucier, Rebecca Reiffenhäuser, and Sahil Singla (2024). \emph{Online Combinatorial Allocations and Auctions with Few Samples}. arXiv:2409.11091v1. \textit{Verified locator:} Sections 4.1 and 4.5; Definitions 4.1--4.2 and constructive proof of Theorem 4.9. \url{https://arxiv.org/html/2409.11091v1}
\par\endgroup

\subsection*{Common conventions: Source mathematical conventions}

\textbf{Finite games.} For a finite set $A$, $\Delta(A)$ is its probability
simplex. Players choose independent mixed strategies in Nash equilibrium;
correlation is allowed only when explicitly part of the solution concept.
In a game with payoff functions $u_i$ and strategy profile $x$, an additive
$\varepsilon$ Nash equilibrium satisfies
\[
 u_i(x)\geq u_i(x_i',x_{-i})-\varepsilon
 \quad\text{for every player $i$ and every admissible deviation $x_i'$}.
\]
For numerical approximation, payoffs are normalized as locally specified.
An $\varepsilon$ payoff residual is distinct from distance at most
$\varepsilon$ to the set of exact equilibria. Point convergence, convergence
of empirical play, and convergence in distance to an equilibrium set are
also distinct.

\textbf{Computational representations.} Unless stated otherwise, a complexity
question takes rational data with binary encoding; polynomial time is measured
in total bit length. A fixed-accuracy polynomial algorithm is not necessarily
polynomial in $1/\varepsilon$, and neither is necessarily polynomial in
$\log(1/\varepsilon)$. Succinct games and value, demand, or sampling oracles
must declare their interfaces. Exact real solutions require an explicit
representation; an unspecified list of arbitrary real numbers is not a
finite computational output. A claimed algorithmic barrier is meaningful
only for its specified model and reductions.

\textbf{Dynamic games.} Histories, information, observation, and allowable
randomization are part of the model. A stationary strategy depends only on
the current state; a positional strategy in a deterministic graph game
chooses one action at each controlled vertex. Nash equilibrium at an initial
state and equilibrium after every history are different requirements.
An expectation of a pathwise $\liminf$ payoff, a $\liminf$ of expected averages,
a limiting expected average, and a uniform finite-horizon equilibrium are
different criteria. None is silently substituted for another.

\textbf{Allocation and mechanisms.} A complete allocation assigns every item
exactly once unless otherwise stated. Zero-valued goods, negative values,
capacity constraints, feasibility, lotteries, and copies of goods affect
fairness definitions and are specified locally. Dominant-strategy incentive
compatibility, Bayesian incentive compatibility, universal truthfulness,
and truthfulness in expectation impose different inequalities. Whenever
an objective ratio has zero denominator, its convention is stated or those
instances are excluded.

\textbf{Infinite-dimensional models.} Expectations, integrals, stochastic
equations, and optimization objectives require their local regularity,
integrability, filtrations, and admissible domains. A backward--forward
mean-field system is a boundary-value problem; it is not an initial-value
semiflow without an additional construction. Long-time, large-population,
and vanishing-noise limits require an order of limits and a convergence
topology. If a minimum need not be attained, the objective is its infimum
or supremum and the approximation requirement is stated separately.
% END ECONOMICS STATEMENT
\end{document}
